\begin{examplebox}
	\textbf{\textcolor{CExample}{例1（基础）}}

	\textbf{\textcolor{CExample}{题目：}}
	在$\triangle ABC$中，已知$a=4$，$b=5$，$C=60^\circ$，求$\triangle ABC$的面积$S$.

	\textbf{\textcolor{CExample}{解题步骤：}}
	\begin{description}[style=nextline, leftmargin=2.8em, labelsep=0.5em, itemsep=0.45em]
		\item[\textbf{第一步（确认条件）：}] 题目给出了两边$a,b$及其夹角$C$，符合公式$S=\frac12ab\sin C$的使用条件.
			\par\textbf{理由：} 公式要求必须知道两边长度及它们的夹角.
		\item[\textbf{第二步（代入公式）：}]
			\[
				S = \frac12 \cdot 4 \cdot 5 \cdot \sin 60^\circ.
			\]
			\par\textbf{理由：} 直接套用三角形面积公式.
		\item[\textbf{第三步（计算数值）：}]
			\[
				\begin{aligned}
					\sin 60^\circ &= \frac{\sqrt{3}}{2},\\
					S &= \frac12 \cdot 4 \cdot 5 \cdot \frac{\sqrt{3}}{2} \\
					&= 10 \cdot \frac{\sqrt{3}}{2} \\
					&= 5\sqrt{3}.
			\end{aligned}
			\]
			\par\textbf{理由：} 利用特殊角的三角函数值，化简得结果.
	\end{description}
	\textbf{\textcolor{CExample}{关键结论：}}
	$\displaystyle S = 5\sqrt{3}$.

	\medskip
	\textbf{\textcolor{CExample}{例2（稍变形）}}

	\textbf{\textcolor{CExample}{题目：}}
	在$\triangle ABC$中，已知$a=6$，$b=8$，面积$S=12\sqrt{3}$，求角$C$.

	\textbf{\textcolor{CExample}{解题步骤：}}
	\begin{description}[style=nextline, leftmargin=2.8em, labelsep=0.5em, itemsep=0.45em]
		\item[\textbf{第一步（反用公式）：}] 由面积公式$S=\frac12 ab\sin C$，得$12\sqrt{3}=\frac12\cdot6\cdot8\cdot\sin C$.
			\par\textbf{理由：} 已知面积和两边，求夹角正弦.
		\item[\textbf{第二步（求解正弦值）：}] 化简得
			\[
				\begin{aligned}
					12\sqrt{3} &= 24\sin C \\
					\Rightarrow \ \sin C &= \frac{12\sqrt{3}}{24} \\
					&= \frac{\sqrt{3}}{2}.
			\end{aligned}
			\]
			\par\textbf{理由：} 解关于$\sin C$的方程.
		\item[\textbf{第三步（确定角C）：}] 在三角形中，$C\in(0,\pi)$，$\sin C=\frac{\sqrt{3}}{2}$对应的锐角为$60^\circ$，钝角为$120^\circ$，故 $C=60^\circ$ 或 $120^\circ$.
			\par\textbf{理由：} 正弦值在$(0,\pi)$内对应两个可能角，需根据三角形内角和判断取舍（本题暂不排除）.
	\end{description}
	\textbf{\textcolor{CExample}{关键结论：}}
	$C=60^\circ$或$C=120^\circ$.

	\medskip
	\textbf{\textcolor{CExample}{例3（综合应用）}}

	\textbf{\textcolor{CExample}{题目：}}
	在$\triangle ABC$中，已知$a=7$，$b=5$，$\cos C = -\frac17$，求$\triangle ABC$的面积$S$.

	\textbf{\textcolor{CExample}{解题步骤：}}
	\begin{description}[style=nextline, leftmargin=2.8em, labelsep=0.5em, itemsep=0.45em]
		\item[\textbf{第一步（求夹角正弦）：}] 由同角三角函数关系$\sin^2 C + \cos^2 C = 1$，得
			\[
				\begin{aligned}
					\sin C &= \sqrt{1 - \left(-\frac17\right)^2} \\
					&= \sqrt{1 - \frac{1}{49}} \\
					&= \sqrt{\frac{48}{49}} \\
					&= \frac{4\sqrt{3}}{7}.
			\end{aligned}
			\]
			又因为$C\in(0,\pi)$，$\cos C<0$说明$C$为钝角，但$\sin C>0$仍成立.
			\par\textbf{理由：} 利用平方关系求正弦，需注意符号（三角形内角正弦恒正）.
		\item[\textbf{第二步（代入面积公式）：}]
			\[
				\begin{aligned}
					S &= \frac12 ab\sin C \\
					&= \frac12 \cdot 7 \cdot 5 \cdot \frac{4\sqrt{3}}{7}.
			\end{aligned}
			\]
			\par\textbf{理由：} 已知两边$a,b$及其夹角$C$，直接使用公式.
		\item[\textbf{第三步（化简得结果）：}]
			\[
				\begin{aligned}
					S &= \frac12 \cdot 5 \cdot 4\sqrt{3} \\
					&= 10\sqrt{3}.
			\end{aligned}
			\]
			\par\textbf{理由：} 约去公因数7，简化计算.
	\end{description}
	\textbf{\textcolor{CExample}{关键结论：}}
	$S=10\sqrt{3}$.
\end{examplebox}
